\documentclass[11pt]{article}
\usepackage[a4paper,margin=1.9cm,top=2.4cm,bottom=2.2cm]{geometry}
\usepackage{amsmath,amssymb,amsfonts,fontspec,microtype,fancyhdr,lastpage,enumitem,needspace}
\usepackage[most]{tcolorbox}
\usepackage{xcolor}
\definecolor{ink}{HTML}{1F3A5F}\definecolor{soft}{HTML}{EEF3F9}\definecolor{ans}{HTML}{F3F8F1}\definecolor{ansb}{HTML}{3C7A3C}
\pagestyle{fancy}\fancyhf{}
\lhead{\small\color{ink}\textbf{Linear Algebra I} \textperiodcentered\ Week 3 --- Span, Linear Independence \& Setting Up Linear Systems}
\rhead{}\cfoot{\small Mr Malek Thiab \textperiodcentered\ Worksheet with Worked Answers \textperiodcentered\ Page \thepage\ of \pageref{LastPage}}
\renewcommand{\headrulewidth}{0.4pt}
\setlength{\parindent}{0pt}\setlength{\parskip}{4pt}
\newcommand{\lv}[1]{\textsc{\small #1}}
\begin{document}
{\Large\bfseries\color{ink} Linear Algebra I --- Week 3}\\[2pt]
{\large\bfseries\color{ink} Span, Linear Independence \& Setting Up Linear Systems}\\[3pt]
{\small Name: \underline{\hspace{5cm}}\hfill Date: \underline{\hspace{3cm}}\qquad All work \textbf{by hand} --- no calculator, no software.}
\vspace{4pt}\hrule\vspace{8pt}

\begin{tcolorbox}[colback=soft,colframe=ink,title={\bfseries Key Facts}]
\textbf{Linear combination and span.} A linear combination of $v_1,\dots,v_k\in V$ is $a_1v_1+\dots+a_kv_k$ with $a_i\in\mathbb{F}$. $\operatorname{span}\{v_1,\dots,v_k\}$ is the set of all of them; $\operatorname{span}\varnothing=\{\mathbf 0\}$. It is a subspace, and the smallest subspace containing $v_1,\dots,v_k$.

\textbf{Setting up the system.} To test $w\in\operatorname{span}\{v_1,\dots,v_k\}$ in $\mathbb{F}^{n}$, solve $a_1v_1+\dots+a_kv_k=w$. Coordinatewise this is $n$ equations in $k$ unknowns, with augmented matrix whose columns are $v_1,\dots,v_k\mid w$. Solution exists $\iff w\in$ span.

\textbf{Linear independence.} $v_1,\dots,v_k$ are \emph{independent} if $a_1v_1+\dots+a_kv_k=\mathbf 0$ forces $a_1=\dots=a_k=0$; otherwise \emph{dependent}. Test: solve the homogeneous system; only $\mathbf 0$ $\Rightarrow$ independent, a nonzero solution $\Rightarrow$ dependent (and it gives a relation).

\textbf{Useful facts.} A set containing $\mathbf 0$ is dependent. One vector is independent iff it is nonzero. Two vectors are independent iff neither is a scalar multiple of the other. $v_1,\dots,v_k$ dependent $\iff$ some $v_j\in\operatorname{span}$ of the others. If $w\in\operatorname{span}\{v_1,\dots,v_k\}$, adding $w$ does not enlarge the span. Independence and span depend on the field: $\{(1,i),(i,-1)\}$ is dependent over $\mathbb{C}$ but independent over $\mathbb{R}$.

\textbf{Counting.} More homogeneous unknowns than equations forces a nonzero solution; so more than $n$ vectors in $\mathbb{F}^{n}$ are dependent.
\end{tcolorbox}

\newpage
\section*{\large\color{ink} Worked Examples}
\begin{tcolorbox}[colback=white,colframe=ink!60,title={\bfseries Example 1 \textperiodcentered\ Membership in a span}]Is $(4,7)\in\operatorname{span}\{(1,2),(3,5)\}$ in $\mathbb{R}^{2}$?
\par\smallskip
\textbf{Solution.} \begin{enumerate}[leftmargin=1.6em,itemsep=1pt,topsep=2pt,label=\arabic*.]\item Solve $a(1,2)+b(3,5)=(4,7)$: $a+3b=4$, $2a+5b=7$.
\item From the first, $a=4-3b$. Substitute: $8-6b+5b=7$, so $b=1$, $a=1$.
\item Check: $(1,2)+(3,5)=(4,7)$.
\end{enumerate}\textbf{Answer:} Yes: $(4,7)=1(1,2)+1(3,5)$.
\end{tcolorbox}
\begin{tcolorbox}[colback=white,colframe=ink!60,title={\bfseries Example 2 \textperiodcentered\ Membership (positive)}]Is $(1,2,3)\in\operatorname{span}\{(1,0,1),(0,1,1)\}$?
\par\smallskip
\textbf{Solution.} \begin{enumerate}[leftmargin=1.6em,itemsep=1pt,topsep=2pt,label=\arabic*.]\item $a(1,0,1)+b(0,1,1)=(a,\ b,\ a+b)$.
\item Match: $a=1$, $b=2$; third coordinate $a+b=3$ \checkmark.
\end{enumerate}\textbf{Answer:} Yes: $(1,2,3)=1(1,0,1)+2(0,1,1)$.
\end{tcolorbox}
\begin{tcolorbox}[colback=white,colframe=ink!60,title={\bfseries Example 3 \textperiodcentered\ Membership (negative)}]Is $(1,2,4)\in\operatorname{span}\{(1,0,1),(0,1,1)\}$?
\par\smallskip
\textbf{Solution.} \begin{enumerate}[leftmargin=1.6em,itemsep=1pt,topsep=2pt,label=\arabic*.]\item As before the combination is $(a,b,a+b)$.
\item Match: $a=1$, $b=2$, but then $a+b=3\neq4$.
\end{enumerate}\textbf{Answer:} No.
\end{tcolorbox}
\begin{tcolorbox}[colback=white,colframe=ink!60,title={\bfseries Example 4 \textperiodcentered\ Dependence, finding a relation}]Test $\{(1,2,1),(2,1,0),(1,-1,-1)\}$ for independence.
\par\smallskip
\textbf{Solution.} \begin{enumerate}[leftmargin=1.6em,itemsep=1pt,topsep=2pt,label=\arabic*.]\item Solve $a(1,2,1)+b(2,1,0)+c(1,-1,-1)=\mathbf 0$: $a+2b+c=0,\ 2a+b-c=0,\ a-c=0$.
\item Third: $c=a$. First: $2a+2b=0$, so $b=-a$. Second: $2a-a-a=0$ \checkmark (automatically).
\item So $(a,b,c)=(1,-1,1)$ is a nontrivial solution. Check: $v_1-v_2+v_3=(1-2+1,\ 2-1-1,\ 1-0-1)=\mathbf 0$.
\end{enumerate}\textbf{Answer:} Linearly dependent: $v_1-v_2+v_3=\mathbf 0$.
\end{tcolorbox}
\begin{tcolorbox}[colback=white,colframe=ink!60,title={\bfseries Example 5 \textperiodcentered\ Independence in $P_{2}$}]Show $\{1+x,\ x+x^{2},\ 1+x^{2}\}$ is linearly independent in $P_{2}(\mathbb{R})$.
\par\smallskip
\textbf{Solution.} \begin{enumerate}[leftmargin=1.6em,itemsep=1pt,topsep=2pt,label=\arabic*.]\item Suppose $a(1+x)+b(x+x^{2})+c(1+x^{2})=\mathbf 0$. Compare coefficients.
\item Constant: $a+c=0$. $x$: $a+b=0$. $x^{2}$: $b+c=0$.
\item So $a=-c$ and $b=-a=c$; then $b+c=2c=0$, so $c=0$, hence $a=b=0$.
\end{enumerate}\textbf{Answer:} Linearly independent.
\end{tcolorbox}
\begin{tcolorbox}[colback=white,colframe=ink!60,title={\bfseries Example 6 \textperiodcentered\ The field matters}]Is $\{(1,i),(i,-1)\}$ linearly independent in $\mathbb{C}^{2}$ over $\mathbb{C}$? Over $\mathbb{R}$?
\par\smallskip
\textbf{Solution.} \begin{enumerate}[leftmargin=1.6em,itemsep=1pt,topsep=2pt,label=\arabic*.]\item Note $i\cdot(1,i)=(i,\,i^{2})=(i,-1)$.
\item Over $\mathbb{C}$: $i\,v_1-v_2=\mathbf 0$, a nontrivial relation, so dependent.
\item Over $\mathbb{R}$: $a(1,i)+b(i,-1)=\mathbf 0$ with real $a,b$ gives first coordinate $a+bi=0$, so $a=b=0$: independent.
\end{enumerate}\textbf{Answer:} Dependent over $\mathbb{C}$; independent over $\mathbb{R}$.
\end{tcolorbox}
\begin{tcolorbox}[colback=white,colframe=ink!60,title={\bfseries Example 7 \textperiodcentered\ Proof: dependence means one vector is redundant}]Prove: if $v_1,\dots,v_k$ are linearly dependent, then some $v_j$ is a linear combination of the others.
\par\smallskip
\textbf{Solution.} \begin{enumerate}[leftmargin=1.6em,itemsep=1pt,topsep=2pt,label=\arabic*.]\item Take a nontrivial relation $a_1v_1+\dots+a_kv_k=\mathbf 0$ and pick $j$ with $a_j\neq0$.
\item Isolate: $a_jv_j=-\sum_{i\neq j}a_iv_i$.
\item Divide by $a_j$ (allowed because $a_j\ne0$ in a field): $v_j=-\sum_{i\ne j}\dfrac{a_i}{a_j}v_i$.
\end{enumerate}\textbf{Answer:} $v_j\in\operatorname{span}\{v_i:i\neq j\}$.
\end{tcolorbox}
\begin{tcolorbox}[colback=white,colframe=ink!60,title={\bfseries Example 8 \textperiodcentered\ Describing a span by an equation}]Find the condition on $(a,b,c)$ for $(a,b,c)\in\operatorname{span}\{(1,1,0),(0,1,1)\}$.
\par\smallskip
\textbf{Solution.} \begin{enumerate}[leftmargin=1.6em,itemsep=1pt,topsep=2pt,label=\arabic*.]\item $x(1,1,0)+y(0,1,1)=(x,\ x+y,\ y)$.
\item Matching: $x=a$, $y=c$, and $b=x+y=a+c$.
\end{enumerate}\textbf{Answer:} $a-b+c=0$
\end{tcolorbox}
\newpage
\section*{\large\color{ink} Practice Problems}
\Needspace{10\baselineskip}\subsection*{\normalsize\color{ink} Part A --- Linear Combinations and Span}
\Needspace{9\baselineskip}\begin{minipage}{\linewidth}\textbf{1.} \hfill{\scriptsize\color{ink}[Foundational]}\par\nopagebreak Write $(5,-1)$ as a linear combination of $(1,1)$ and $(1,-1)$ in $\mathbb{R}^{2}$.
\end{minipage}\par\nopagebreak
\begin{tcolorbox}[breakable,colback=ans,colframe=ansb,boxrule=0.5pt,left=4pt,right=4pt,top=2pt,bottom=2pt]\begin{enumerate}[leftmargin=2.2em,itemsep=1pt,topsep=1pt,label=\textbf{Step \arabic*.}]\item Set $a(1,1)+b(1,-1)=(5,-1)$: $a+b=5$ and $a-b=-1$.
\item Add the equations: $2a=4$, so $a=2$. Then $b=5-2=3$.
\item Check: $2(1,1)+3(1,-1)=(2+3,\,2-3)=(5,-1)$.
\end{enumerate}\textbf{Answer:} $(5,-1)=2(1,1)+3(1,-1)$\end{tcolorbox}
\par\vspace{2pt}
\Needspace{9\baselineskip}\begin{minipage}{\linewidth}\textbf{2.} \hfill{\scriptsize\color{ink}[Foundational]}\par\nopagebreak Decide whether $(2,3,4)\in\operatorname{span}\{(1,1,1),(0,1,2)\}$ in $\mathbb{R}^{3}$. If so, give the coefficients.
\end{minipage}\par\nopagebreak
\begin{tcolorbox}[breakable,colback=ans,colframe=ansb,boxrule=0.5pt,left=4pt,right=4pt,top=2pt,bottom=2pt]\begin{enumerate}[leftmargin=2.2em,itemsep=1pt,topsep=1pt,label=\textbf{Step \arabic*.}]\item $a(1,1,1)+b(0,1,2)=(a,\ a+b,\ a+2b)$. Set equal to $(2,3,4)$.
\item First coordinate: $a=2$. Second: $2+b=3$, so $b=1$.
\item Third coordinate check: $a+2b=2+2=4$ \checkmark.
\end{enumerate}\textbf{Answer:} Yes: $(2,3,4)=2(1,1,1)+1\,(0,1,2)$.\end{tcolorbox}
\par\vspace{2pt}
\Needspace{9\baselineskip}\begin{minipage}{\linewidth}\textbf{3.} \hfill{\scriptsize\color{ink}[Foundational]}\par\nopagebreak Decide whether $(1,0,0)\in\operatorname{span}\{(1,1,0),(0,1,1)\}$ in $\mathbb{R}^{3}$.
\end{minipage}\par\nopagebreak
\begin{tcolorbox}[breakable,colback=ans,colframe=ansb,boxrule=0.5pt,left=4pt,right=4pt,top=2pt,bottom=2pt]\begin{enumerate}[leftmargin=2.2em,itemsep=1pt,topsep=1pt,label=\textbf{Step \arabic*.}]\item $x(1,1,0)+y(0,1,1)=(x,\ x+y,\ y)$. Set equal to $(1,0,0)$.
\item First coordinate: $x=1$. Third: $y=0$.
\item Second coordinate: $x+y=1\neq0$. Contradiction.
\end{enumerate}\textbf{Answer:} No, $(1,0,0)\notin\operatorname{span}\{(1,1,0),(0,1,1)\}$.\end{tcolorbox}
\par\vspace{2pt}
\Needspace{9\baselineskip}\begin{minipage}{\linewidth}\textbf{4.} \hfill{\scriptsize\color{ink}[Intermediate]}\par\nopagebreak Find $k$ such that $(1,2,k)\in\operatorname{span}\{(1,1,0),(1,-1,2)\}$.
\end{minipage}\par\nopagebreak
\begin{tcolorbox}[breakable,colback=ans,colframe=ansb,boxrule=0.5pt,left=4pt,right=4pt,top=2pt,bottom=2pt]\begin{enumerate}[leftmargin=2.2em,itemsep=1pt,topsep=1pt,label=\textbf{Step \arabic*.}]\item $a(1,1,0)+b(1,-1,2)=(a+b,\ a-b,\ 2b)$. Set equal to $(1,2,k)$.
\item $a+b=1$ and $a-b=2$: adding gives $2a=3$, so $a=\tfrac32$; then $b=-\tfrac12$.
\item Third coordinate: $k=2b=-1$.
\item Check: $\tfrac32(1,1,0)-\tfrac12(1,-1,2)=(1,\,2,\,-1)$ \checkmark.
\end{enumerate}\textbf{Answer:} $k=-1$\end{tcolorbox}
\par\vspace{2pt}
\Needspace{9\baselineskip}\begin{minipage}{\linewidth}\textbf{5.} \hfill{\scriptsize\color{ink}[Intermediate]}\par\nopagebreak In $\mathbb{C}^{2}$ over $\mathbb{C}$, write $(1+i,\;2)$ as a linear combination of $(1,i)$ and $(1,-i)$.
\end{minipage}\par\nopagebreak
\begin{tcolorbox}[breakable,colback=ans,colframe=ansb,boxrule=0.5pt,left=4pt,right=4pt,top=2pt,bottom=2pt]\begin{enumerate}[leftmargin=2.2em,itemsep=1pt,topsep=1pt,label=\textbf{Step \arabic*.}]\item Solve $a(1,i)+b(1,-i)=(1+i,2)$: \ $a+b=1+i$ \ and \ $ia-ib=2$.
\item Divide the second equation by $i$: $a-b=\dfrac{2}{i}=-2i$.
\item Add: $2a=1+i-2i=1-i$, so $a=\dfrac{1-i}{2}$. Subtract: $2b=1+i+2i=1+3i$, so $b=\dfrac{1+3i}{2}$.
\item Check: $a+b=\tfrac{2+2i}{2}=1+i$ \checkmark \ and \ $a-b=\tfrac{-4i}{2}=-2i$ \checkmark.
\end{enumerate}\textbf{Answer:} $a=\dfrac{1-i}{2},\ b=\dfrac{1+3i}{2}$\end{tcolorbox}
\par\vspace{2pt}
\Needspace{9\baselineskip}\begin{minipage}{\linewidth}\textbf{6.} \hfill{\scriptsize\color{ink}[Foundational]}\par\nopagebreak Show that $\operatorname{span}\{(1,2),(2,4)\}\subseteq\mathbb{R}^{2}$ is a line, and give its equation.
\end{minipage}\par\nopagebreak
\begin{tcolorbox}[breakable,colback=ans,colframe=ansb,boxrule=0.5pt,left=4pt,right=4pt,top=2pt,bottom=2pt]\begin{enumerate}[leftmargin=2.2em,itemsep=1pt,topsep=1pt,label=\textbf{Step \arabic*.}]\item A general element is $a(1,2)+b(2,4)=(a+2b,\ 2a+4b)=(a+2b)(1,2)$.
\item Put $t=a+2b$; as $a,b$ vary, $t$ takes every real value.
\item So the span is $\{(t,2t):t\in\mathbb{R}\}$, i.e.\ all points with $y=2x$.
\end{enumerate}\textbf{Answer:} The line $y=2x$ through the origin.\end{tcolorbox}
\par\vspace{2pt}
\Needspace{9\baselineskip}\begin{minipage}{\linewidth}\textbf{7.} \hfill{\scriptsize\color{ink}[Intermediate]}\par\nopagebreak Describe $\operatorname{span}\{(1,0,2),(0,1,-1)\}\subseteq\mathbb{R}^{3}$ by a single linear equation in $x,y,z$.
\end{minipage}\par\nopagebreak
\begin{tcolorbox}[breakable,colback=ans,colframe=ansb,boxrule=0.5pt,left=4pt,right=4pt,top=2pt,bottom=2pt]\begin{enumerate}[leftmargin=2.2em,itemsep=1pt,topsep=1pt,label=\textbf{Step \arabic*.}]\item $a(1,0,2)+b(0,1,-1)=(a,\ b,\ 2a-b)$.
\item So $x=a$, $y=b$, $z=2a-b=2x-y$.
\item Rearrange: $2x-y-z=0$. Conversely every $(x,y,2x-y)$ is obtained with $a=x,\ b=y$.
\end{enumerate}\textbf{Answer:} $\{(x,y,z):2x-y-z=0\}$\end{tcolorbox}
\par\vspace{2pt}
\Needspace{9\baselineskip}\begin{minipage}{\linewidth}\textbf{8.} \hfill{\scriptsize\color{ink}[Challenge]}\par\nopagebreak Prove that for any vectors $v_1,\dots,v_k$ in a vector space $V$, the set $\operatorname{span}\{v_1,\dots,v_k\}$ is a subspace of $V$.
\end{minipage}\par\nopagebreak
\begin{tcolorbox}[breakable,colback=ans,colframe=ansb,boxrule=0.5pt,left=4pt,right=4pt,top=2pt,bottom=2pt]\begin{enumerate}[leftmargin=2.2em,itemsep=1pt,topsep=1pt,label=\textbf{Step \arabic*.}]\item Zero vector: $\mathbf 0=0v_1+\dots+0v_k$, so $\mathbf 0\in\operatorname{span}$.
\item Addition: $\sum a_iv_i+\sum b_iv_i=\sum(a_i+b_i)v_i$, again a linear combination.
\item Scalars: $\lambda\sum a_iv_i=\sum(\lambda a_i)v_i$, again a linear combination.
\item All three subspace conditions hold.
\end{enumerate}\textbf{Answer:} $\operatorname{span}\{v_1,\dots,v_k\}$ is a subspace of $V$.\end{tcolorbox}
\par\vspace{2pt}
\Needspace{10\baselineskip}\subsection*{\normalsize\color{ink} Part B --- Linear Independence and Dependence}
\Needspace{9\baselineskip}\begin{minipage}{\linewidth}\textbf{9.} \hfill{\scriptsize\color{ink}[Foundational]}\par\nopagebreak Is $\{(1,2),(3,6)\}$ linearly independent in $\mathbb{R}^{2}$?
\end{minipage}\par\nopagebreak
\begin{tcolorbox}[breakable,colback=ans,colframe=ansb,boxrule=0.5pt,left=4pt,right=4pt,top=2pt,bottom=2pt]\begin{enumerate}[leftmargin=2.2em,itemsep=1pt,topsep=1pt,label=\textbf{Step \arabic*.}]\item Observe $(3,6)=3(1,2)$.
\item Hence $3(1,2)-1\,(3,6)=\mathbf 0$ is a nontrivial relation.
\end{enumerate}\textbf{Answer:} Linearly dependent.\end{tcolorbox}
\par\vspace{2pt}
\Needspace{9\baselineskip}\begin{minipage}{\linewidth}\textbf{10.} \hfill{\scriptsize\color{ink}[Foundational]}\par\nopagebreak Is $\{(1,2),(2,3)\}$ linearly independent in $\mathbb{R}^{2}$?
\end{minipage}\par\nopagebreak
\begin{tcolorbox}[breakable,colback=ans,colframe=ansb,boxrule=0.5pt,left=4pt,right=4pt,top=2pt,bottom=2pt]\begin{enumerate}[leftmargin=2.2em,itemsep=1pt,topsep=1pt,label=\textbf{Step \arabic*.}]\item Solve $a(1,2)+b(2,3)=\mathbf 0$: $a+2b=0$ and $2a+3b=0$.
\item From the first, $a=-2b$. Substitute: $-4b+3b=-b=0$, so $b=0$ and $a=0$.
\item Only the trivial solution.
\end{enumerate}\textbf{Answer:} Linearly independent.\end{tcolorbox}
\par\vspace{2pt}
\Needspace{9\baselineskip}\begin{minipage}{\linewidth}\textbf{11.} \hfill{\scriptsize\color{ink}[Intermediate]}\par\nopagebreak Is $\{(1,0,1),(1,1,0),(0,1,1)\}$ linearly independent in $\mathbb{R}^{3}$?
\end{minipage}\par\nopagebreak
\begin{tcolorbox}[breakable,colback=ans,colframe=ansb,boxrule=0.5pt,left=4pt,right=4pt,top=2pt,bottom=2pt]\begin{enumerate}[leftmargin=2.2em,itemsep=1pt,topsep=1pt,label=\textbf{Step \arabic*.}]\item Solve $a(1,0,1)+b(1,1,0)+c(0,1,1)=\mathbf 0$: \ $a+b=0,\ \ b+c=0,\ \ a+c=0$.
\item So $a=-b$ and $c=-b$. Then $a+c=-2b=0$, giving $b=0$.
\item Hence $a=b=c=0$.
\end{enumerate}\textbf{Answer:} Linearly independent.\end{tcolorbox}
\par\vspace{2pt}
\Needspace{9\baselineskip}\begin{minipage}{\linewidth}\textbf{12.} \hfill{\scriptsize\color{ink}[Intermediate]}\par\nopagebreak Is $\{(1,2,3),(4,5,6),(7,8,9)\}$ linearly independent in $\mathbb{R}^{3}$? If not, give a nontrivial relation.
\end{minipage}\par\nopagebreak
\begin{tcolorbox}[breakable,colback=ans,colframe=ansb,boxrule=0.5pt,left=4pt,right=4pt,top=2pt,bottom=2pt]\begin{enumerate}[leftmargin=2.2em,itemsep=1pt,topsep=1pt,label=\textbf{Step \arabic*.}]\item Solve $a v_1+bv_2+cv_3=\mathbf 0$: $a+4b+7c=0,\ 2a+5b+8c=0,\ 3a+6b+9c=0$.
\item Eq.\,2 $-$ Eq.\,1: $a+b+c=0$. Eq.\,3 $-$ Eq.\,2: $a+b+c=0$ (same equation). So the system has a free variable.
\item Eq.\,1 $-$ $(a+b+c)$: $3b+6c=0$, so $b=-2c$; then $a=-b-c=c$. Take $c=1$: $(a,b,c)=(1,-2,1)$.
\item Check: $v_1-2v_2+v_3=(1-8+7,\ 2-10+8,\ 3-12+9)=(0,0,0)$ \checkmark.
\end{enumerate}\textbf{Answer:} Linearly dependent: $v_1-2v_2+v_3=\mathbf 0$.\end{tcolorbox}
\par\vspace{2pt}
\Needspace{9\baselineskip}\begin{minipage}{\linewidth}\textbf{13.} \hfill{\scriptsize\color{ink}[Challenge]}\par\nopagebreak Find all real $k$ for which $\{(1,1,1),(1,2,3),(1,3,k)\}$ is linearly dependent.
\end{minipage}\par\nopagebreak
\begin{tcolorbox}[breakable,colback=ans,colframe=ansb,boxrule=0.5pt,left=4pt,right=4pt,top=2pt,bottom=2pt]\begin{enumerate}[leftmargin=2.2em,itemsep=1pt,topsep=1pt,label=\textbf{Step \arabic*.}]\item Solve $a(1,1,1)+b(1,2,3)+c(1,3,k)=\mathbf 0$: $a+b+c=0,\ \ a+2b+3c=0,\ \ a+3b+kc=0$.
\item Eq.\,2 $-$ Eq.\,1: $b+2c=0$. Eq.\,3 $-$ Eq.\,2: $b+(k-3)c=0$.
\item Subtract these two: $(k-5)c=0$.
\item If $k\ne5$: $c=0$, then $b=0,\ a=0$ (independent). If $k=5$: $c$ is free, e.g.\ $c=1,\ b=-2,\ a=1$ (dependent).
\end{enumerate}\textbf{Answer:} Dependent exactly when $k=5$.\end{tcolorbox}
\par\vspace{2pt}
\Needspace{9\baselineskip}\begin{minipage}{\linewidth}\textbf{14.} \hfill{\scriptsize\color{ink}[Foundational]}\par\nopagebreak In $P_{2}(\mathbb{R})$, is $\{1+x,\ 1-x,\ 2\}$ linearly independent?
\end{minipage}\par\nopagebreak
\begin{tcolorbox}[breakable,colback=ans,colframe=ansb,boxrule=0.5pt,left=4pt,right=4pt,top=2pt,bottom=2pt]\begin{enumerate}[leftmargin=2.2em,itemsep=1pt,topsep=1pt,label=\textbf{Step \arabic*.}]\item Add the first two: $(1+x)+(1-x)=2$.
\item So $(1+x)+(1-x)-1\cdot2=\mathbf 0$ is a nontrivial relation.
\end{enumerate}\textbf{Answer:} Linearly dependent.\end{tcolorbox}
\par\vspace{2pt}
\Needspace{9\baselineskip}\begin{minipage}{\linewidth}\textbf{15.} \hfill{\scriptsize\color{ink}[Intermediate]}\par\nopagebreak In $\mathbb{C}^{2}$, consider $\{(1,i),\ (1+i,\ i-1)\}$. (a) Is it linearly independent over $\mathbb{C}$? (b) Is it linearly independent over $\mathbb{R}$?
\end{minipage}\par\nopagebreak
\begin{tcolorbox}[breakable,colback=ans,colframe=ansb,boxrule=0.5pt,left=4pt,right=4pt,top=2pt,bottom=2pt]\begin{enumerate}[leftmargin=2.2em,itemsep=1pt,topsep=1pt,label=\textbf{Step \arabic*.}]\item (a) Compute $(1+i)(1,i)=(1+i,\ i+i^{2})=(1+i,\ i-1)$, which is the second vector.
\item So $(1+i)\,v_1-v_2=\mathbf 0$: dependent over $\mathbb{C}$.
\item (b) Take real $a,b$ with $a(1,i)+b(1+i,\,i-1)=\mathbf 0$. First coordinate: $a+b+bi=0$.
\item Imaginary part gives $b=0$; then real part gives $a=0$. Only the trivial solution over $\mathbb{R}$.
\end{enumerate}\textbf{Answer:} (a) dependent over $\mathbb{C}$; (b) independent over $\mathbb{R}$.\end{tcolorbox}
\par\vspace{2pt}
\Needspace{9\baselineskip}\begin{minipage}{\linewidth}\textbf{16.} \hfill{\scriptsize\color{ink}[Challenge]}\par\nopagebreak Let $u,v,w$ be linearly independent in a real vector space. Prove that $u+v,\ v+w,\ u+w$ are linearly independent.
\end{minipage}\par\nopagebreak
\begin{tcolorbox}[breakable,colback=ans,colframe=ansb,boxrule=0.5pt,left=4pt,right=4pt,top=2pt,bottom=2pt]\begin{enumerate}[leftmargin=2.2em,itemsep=1pt,topsep=1pt,label=\textbf{Step \arabic*.}]\item Suppose $a(u+v)+b(v+w)+c(u+w)=\mathbf 0$.
\item Regroup: $(a+c)u+(a+b)v+(b+c)w=\mathbf 0$.
\item By independence of $u,v,w$: $a+c=0,\ a+b=0,\ b+c=0$.
\item Adding: $2(a+b+c)=0$, so $a+b+c=0$. Then $a=(a+b+c)-(b+c)=0$, and similarly $b=0$, $c=0$.
\end{enumerate}\textbf{Answer:} Only the trivial relation exists, so they are linearly independent.\end{tcolorbox}
\par\vspace{2pt}
\Needspace{10\baselineskip}\subsection*{\normalsize\color{ink} Part C --- Setting Up Systems, Spanning Sets and Proofs}
\Needspace{9\baselineskip}\begin{minipage}{\linewidth}\textbf{17.} \hfill{\scriptsize\color{ink}[Foundational]}\par\nopagebreak Set up and solve the system that decides whether $(1,2,3)\in\operatorname{span}\{(1,0,1),(2,1,0),(0,1,1)\}$.
\end{minipage}\par\nopagebreak
\begin{tcolorbox}[breakable,colback=ans,colframe=ansb,boxrule=0.5pt,left=4pt,right=4pt,top=2pt,bottom=2pt]\begin{enumerate}[leftmargin=2.2em,itemsep=1pt,topsep=1pt,label=\textbf{Step \arabic*.}]\item Seek $a,b,c$ with $a(1,0,1)+b(2,1,0)+c(0,1,1)=(1,2,3)$. Coordinatewise: $a+2b=1,\ \ b+c=2,\ \ a+c=3$.
\item Augmented matrix (columns are the three vectors, last column is the target): $\left(\begin{array}{ccc|c}1&2&0&1\\0&1&1&2\\1&0&1&3\end{array}\right)$.
\item From Eq.\,1: $a=1-2b$. From Eq.\,2: $c=2-b$. Substitute into Eq.\,3: $3-3b=3$, so $b=0$, $a=1$, $c=2$.
\item Check: $1(1,0,1)+0(2,1,0)+2(0,1,1)=(1,2,3)$ \checkmark.
\end{enumerate}\textbf{Answer:} Yes: $(1,2,3)=1(1,0,1)+0(2,1,0)+2(0,1,1)$.\end{tcolorbox}
\par\vspace{2pt}
\Needspace{9\baselineskip}\begin{minipage}{\linewidth}\textbf{18.} \hfill{\scriptsize\color{ink}[Intermediate]}\par\nopagebreak Find the condition on $(a,b,c)$ for $(a,b,c)\in\operatorname{span}\{(1,2,0),(0,1,1)\}$.
\end{minipage}\par\nopagebreak
\begin{tcolorbox}[breakable,colback=ans,colframe=ansb,boxrule=0.5pt,left=4pt,right=4pt,top=2pt,bottom=2pt]\begin{enumerate}[leftmargin=2.2em,itemsep=1pt,topsep=1pt,label=\textbf{Step \arabic*.}]\item $x(1,2,0)+y(0,1,1)=(x,\ 2x+y,\ y)$.
\item Matching $(a,b,c)$: $x=a$, $y=c$, and $b=2x+y=2a+c$.
\item The condition is $2a-b+c=0$; conversely, if it holds take $x=a,\ y=c$.
\end{enumerate}\textbf{Answer:} $2a-b+c=0$\end{tcolorbox}
\par\vspace{2pt}
\Needspace{9\baselineskip}\begin{minipage}{\linewidth}\textbf{19.} \hfill{\scriptsize\color{ink}[Intermediate]}\par\nopagebreak Find a finite spanning set for $W=\{(x,y,z)\in\mathbb{R}^{3}:x+y+z=0\}$.
\end{minipage}\par\nopagebreak
\begin{tcolorbox}[breakable,colback=ans,colframe=ansb,boxrule=0.5pt,left=4pt,right=4pt,top=2pt,bottom=2pt]\begin{enumerate}[leftmargin=2.2em,itemsep=1pt,topsep=1pt,label=\textbf{Step \arabic*.}]\item Solve for $x$: $x=-y-z$, so every element is $(-y-z,\ y,\ z)$.
\item Split: $(-y-z,y,z)=y(-1,1,0)+z(-1,0,1)$.
\item Hence $W\subseteq\operatorname{span}\{(-1,1,0),(-1,0,1)\}$; both vectors satisfy $x+y+z=0$, so equality holds.
\end{enumerate}\textbf{Answer:} $W=\operatorname{span}\{(-1,1,0),(-1,0,1)\}$\end{tcolorbox}
\par\vspace{2pt}
\Needspace{9\baselineskip}\begin{minipage}{\linewidth}\textbf{20.} \hfill{\scriptsize\color{ink}[Intermediate]}\par\nopagebreak (a) Show that $\{(1,1),(1,2),(2,1)\}$ spans $\mathbb{R}^{2}$. (b) Is it linearly independent?
\end{minipage}\par\nopagebreak
\begin{tcolorbox}[breakable,colback=ans,colframe=ansb,boxrule=0.5pt,left=4pt,right=4pt,top=2pt,bottom=2pt]\begin{enumerate}[leftmargin=2.2em,itemsep=1pt,topsep=1pt,label=\textbf{Step \arabic*.}]\item (a) For a target $(x,y)$ use only the first two: $a+b=x$, $a+2b=y$. Subtract: $b=y-x$, then $a=2x-y$. Every $(x,y)$ is reached.
\item (b) Look for a relation among all three. Try $(2,1)=a(1,1)+b(1,2)$: $a+b=2,\ a+2b=1$ give $b=-1,\ a=3$.
\item Check: $3(1,1)-(1,2)=(2,1)$ \checkmark, so $3(1,1)-(1,2)-(2,1)=\mathbf 0$.
\end{enumerate}\textbf{Answer:} (a) Yes it spans. (b) No, it is linearly dependent.\end{tcolorbox}
\par\vspace{2pt}
\Needspace{9\baselineskip}\begin{minipage}{\linewidth}\textbf{21.} \hfill{\scriptsize\color{ink}[Intermediate]}\par\nopagebreak Show that $\{1,\ 1+x,\ 1+x+x^{2}\}$ spans $P_{2}(\mathbb{R})$.
\end{minipage}\par\nopagebreak
\begin{tcolorbox}[breakable,colback=ans,colframe=ansb,boxrule=0.5pt,left=4pt,right=4pt,top=2pt,bottom=2pt]\begin{enumerate}[leftmargin=2.2em,itemsep=1pt,topsep=1pt,label=\textbf{Step \arabic*.}]\item Let $a+bx+cx^{2}$ be arbitrary. Seek $\alpha,\beta,\gamma$ with $a+bx+cx^{2}=\alpha\cdot1+\beta(1+x)+\gamma(1+x+x^{2})$.
\item Compare coefficients: $x^{2}$: $\gamma=c$. $x$: $\beta+\gamma=b$, so $\beta=b-c$. Constant: $\alpha+\beta+\gamma=a$, so $\alpha=a-b$.
\item Such $\alpha,\beta,\gamma$ exist for every $a,b,c$.
\end{enumerate}\textbf{Answer:} $a+bx+cx^{2}=(a-b)\cdot1+(b-c)(1+x)+c(1+x+x^{2})$; spans.\end{tcolorbox}
\par\vspace{2pt}
\Needspace{9\baselineskip}\begin{minipage}{\linewidth}\textbf{22.} \hfill{\scriptsize\color{ink}[Challenge]}\par\nopagebreak Show that $\left\{\begin{pmatrix}1&0\\0&1\end{pmatrix},\begin{pmatrix}0&1\\1&0\end{pmatrix},\begin{pmatrix}0&1\\-1&0\end{pmatrix},\begin{pmatrix}1&0\\0&-1\end{pmatrix}\right\}$ spans $M_{2\times2}(\mathbb{R})$ and is linearly independent.
\end{minipage}\par\nopagebreak
\begin{tcolorbox}[breakable,colback=ans,colframe=ansb,boxrule=0.5pt,left=4pt,right=4pt,top=2pt,bottom=2pt]\begin{enumerate}[leftmargin=2.2em,itemsep=1pt,topsep=1pt,label=\textbf{Step \arabic*.}]\item Let $\begin{pmatrix}p&q\\r&s\end{pmatrix}=\alpha I+\beta\begin{pmatrix}0&1\\1&0\end{pmatrix}+\gamma\begin{pmatrix}0&1\\-1&0\end{pmatrix}+\delta\begin{pmatrix}1&0\\0&-1\end{pmatrix}$.
\item Entries give: $p=\alpha+\delta$, \ $s=\alpha-\delta$, \ $q=\beta+\gamma$, \ $r=\beta-\gamma$.
\item Solve: $\alpha=\tfrac{p+s}{2},\ \delta=\tfrac{p-s}{2},\ \beta=\tfrac{q+r}{2},\ \gamma=\tfrac{q-r}{2}$. So every matrix is reached: spanning.
\item Independence: if the matrix is $\mathbf 0$ then $p=q=r=s=0$, so the formulas give $\alpha=\beta=\gamma=\delta=0$.
\end{enumerate}\textbf{Answer:} The set spans $M_{2\times2}(\mathbb{R})$ and is linearly independent.\end{tcolorbox}
\par\vspace{2pt}
\Needspace{9\baselineskip}\begin{minipage}{\linewidth}\textbf{23.} \hfill{\scriptsize\color{ink}[Challenge]}\par\nopagebreak Prove: if $w\in\operatorname{span}\{v_1,\dots,v_k\}$ then $\operatorname{span}\{v_1,\dots,v_k,w\}=\operatorname{span}\{v_1,\dots,v_k\}$.
\end{minipage}\par\nopagebreak
\begin{tcolorbox}[breakable,colback=ans,colframe=ansb,boxrule=0.5pt,left=4pt,right=4pt,top=2pt,bottom=2pt]\begin{enumerate}[leftmargin=2.2em,itemsep=1pt,topsep=1pt,label=\textbf{Step \arabic*.}]\item ($\supseteq$) Any combination of $v_1,\dots,v_k$ is a combination of $v_1,\dots,v_k,w$ with coefficient $0$ on $w$.
\item Write $w=\sum c_iv_i$. ($\subseteq$) Take $\sum a_iv_i+bw$.
\item Substitute: $\sum a_iv_i+b\sum c_iv_i=\sum(a_i+bc_i)v_i$, a combination of $v_1,\dots,v_k$.
\item Both inclusions hold.
\end{enumerate}\textbf{Answer:} The two spans are equal.\end{tcolorbox}
\par\vspace{2pt}
\Needspace{9\baselineskip}\begin{minipage}{\linewidth}\textbf{24.} \hfill{\scriptsize\color{ink}[Challenge]}\par\nopagebreak Let $v_1,\dots,v_k$ be linearly independent and $v_{k+1}\notin\operatorname{span}\{v_1,\dots,v_k\}$. Prove that $v_1,\dots,v_{k+1}$ are linearly independent.
\end{minipage}\par\nopagebreak
\begin{tcolorbox}[breakable,colback=ans,colframe=ansb,boxrule=0.5pt,left=4pt,right=4pt,top=2pt,bottom=2pt]\begin{enumerate}[leftmargin=2.2em,itemsep=1pt,topsep=1pt,label=\textbf{Step \arabic*.}]\item Suppose $a_1v_1+\dots+a_kv_k+a_{k+1}v_{k+1}=\mathbf 0$.
\item If $a_{k+1}\neq0$: $v_{k+1}=-\sum_{i\le k}\dfrac{a_i}{a_{k+1}}v_i\in\operatorname{span}\{v_1,\dots,v_k\}$, a contradiction.
\item So $a_{k+1}=0$, leaving $a_1v_1+\dots+a_kv_k=\mathbf 0$.
\item By independence of $v_1,\dots,v_k$: $a_1=\dots=a_k=0$.
\end{enumerate}\textbf{Answer:} All coefficients vanish, so $v_1,\dots,v_{k+1}$ are linearly independent.\end{tcolorbox}
\par\vspace{2pt}
\end{document}
